\section{The original affine native differential in its original section carrier}

\begin{lemma}[Original affine section presentations]
\label{lem:kn-canonical-affine-native-readers}
\lean{EnvelopeProperty.Vanishing.NadelGlobalCanonicalTrace.restrictedNativeSection,EnvelopeProperty.Vanishing.NadelGlobalCanonicalTrace.restrictedModuleSpecGammaToRaw,EnvelopeProperty.Vanishing.NadelGlobalCanonicalTrace.restrictedModuleSpecGammaFromNative,EnvelopeProperty.Vanishing.NadelGlobalCanonicalTrace.restrictedModuleSpecGammaToRaw_original,EnvelopeProperty.Vanishing.NadelGlobalCanonicalTrace.affinePushforwardRawSection,EnvelopeProperty.Vanishing.NadelGlobalCanonicalTrace.affinePushforwardRawRead,EnvelopeProperty.Vanishing.NadelGlobalCanonicalTrace.affinePushforwardRawRead_section,EnvelopeProperty.Vanishing.NadelGlobalCanonicalTrace.affinePushforwardRawSection_read}
\leanok
Let $R\to S$ be a homomorphism of commutative rings, and let $M$ be
an $\mathcal O_{\operatorname{Spec}S}$-module. The native affine section
module, its restriction of scalars to $R$, and the raw global-section
carrier of $M$ have their canonical original underlying additive groups.
The maps that forget the restricted scalar presentation or read a native
section as its raw section preserve the section itself. Reading a top
pushforward section of $M$ likewise preserves the same section, and is
inverse to its canonical presentation as a pushforward section.
\end{lemma}
\begin{proof}
Restriction of scalars changes the scalar action and keeps the underlying
additive group. The native affine global-section functor obtains the
$S$-action from the whole structure-section action through the canonical
spectrum section isomorphism. Its raw reader keeps the underlying section.
Finally, for $\beta:\operatorname{Spec}S\to\operatorname{Spec}R$,
the inverse image of the whole target is the whole source. The top
pushforward section and the source section are therefore the same
section. Each asserted readback or inverse identity follows directly
from these definitions; no new section or coefficient is chosen.
\end{proof}

\begin{lemma}[The actual map on the original raw carrier]
\label{lem:kn-canonical-affine-native-raw-map}
\lean{EnvelopeProperty.Vanishing.NadelGlobalCanonicalTrace.affineRawModuleMap,EnvelopeProperty.Vanishing.NadelGlobalCanonicalTrace.affineRawModuleMap_apply,EnvelopeProperty.Vanishing.NadelGlobalCanonicalTrace.moduleSpecGamma_map_toRaw,EnvelopeProperty.Vanishing.NadelGlobalCanonicalTrace.moduleSpecGamma_smul_toRaw,EnvelopeProperty.Vanishing.NadelGlobalCanonicalTrace.affinePushforwardRawSection_map,EnvelopeProperty.Vanishing.NadelGlobalCanonicalTrace.affinePushforwardRawRead_map,EnvelopeProperty.Vanishing.NadelGlobalCanonicalTrace.specPushforwardSectionsIso_hom_raw}
\leanok
For a scheme $X$, $\mathcal O_X$-modules $F,G$, and a morphism $a:F\to G$,
the raw section map is the actual component $a(X)$ on the original raw
section carriers of $F$ and $G$. The native affine section functor reads
as this same component. Scalar multiplication by an affine coefficient
$s\in S$ reads as multiplication by its image in
$\Gamma(\operatorname{Spec}S,\mathcal O)$ under the canonical spectrum
section isomorphism. The canonical affine pushforward comparison also
reads as the original section.
\end{lemma}
\begin{proof}
\uses{lem:kn-canonical-affine-native-readers}
The raw map is defined by evaluating the original sheaf morphism on the
whole open set. The affine section functor applies that same component
and then restricts the scalar action along the spectrum section
isomorphism. This gives both the map identity and the scalar identity.
Pushforward uses the same original component on the inverse-image open;
on whole opens this is the component just described. The canonical
pushforward section comparison therefore gives the same readback.
\end{proof}

\begin{lemma}[The original affine pullback unit]
\label{lem:kn-canonical-affine-native-unit-readback}
\lean{EnvelopeProperty.Vanishing.NadelGlobalCanonicalTrace.affinePullbackUnitTop_eq_unit,EnvelopeProperty.Vanishing.NadelGlobalCanonicalTrace.affinePullbackUnitTop_eq_read_unit}
\leanok
Let $\beta:\operatorname{Spec}S\to\operatorname{Spec}R$ come from the
original ring homomorphism, and let $L$ be an
$\mathcal O_{\operatorname{Spec}R}$-module. The original affine top pullback
unit applied to a section $w$ is exactly the raw source readback of the
actual adjunction unit $L\to\beta_*\beta^*L$ applied to $w$.
\end{lemma}
\begin{proof}
\uses{lem:kn-canonical-affine-native-readers}
The affine top pullback unit is the pullback unit on the whole open,
with the inverse-image whole open identified with the original whole
source. The identity restriction on that whole open fixes every
section. Reading the resulting pushforward section therefore gives the
same original pullback section, as asserted.
\end{proof}

\begin{lemma}[Original field-structure equality transport]
\label{lem:kn-canonical-affine-native-transport}
\lean{EnvelopeProperty.Vanishing.NadelGlobalCanonicalTrace.affineCanonicalGammaTransport_toRaw,EnvelopeProperty.Vanishing.NadelGlobalCanonicalTrace.affineCanonicalGammaTransport_restricted_toRaw}
\leanok
Let $k$ be a field, let $S$ be any commutative ring, and let
$g_1,g_2:\operatorname{Spec}S\to\operatorname{Spec}k$ be the same original
morphism by an equality $h:g_1=g_2$. The original native equality transport
between the affine section modules of $\Omega^n_{g_1}$ and
$\Omega^n_{g_2}$ reads as the actual raw component of the sheaf equality
transport induced by $h$. This remains true after any restriction of
scalars along an original ring homomorphism $R\to S$.
\end{lemma}
\begin{proof}
\uses{lem:kn-canonical-affine-native-readers,lem:kn-canonical-affine-native-raw-map}
The original native transport equals the affine section functor applied
to the actual sheaf equality transport. Replace it by that equal original
map and use the preceding component identity. Alternatively, eliminate
$h$. Both transports then become the identity; the assertion is the
identity readback of the original section. Restriction of scalars does
not change this underlying value.
\end{proof}

\begin{theorem}[Actual native differential readback over arbitrary affine rings]
\label{thm:kn-canonical-affine-native-differential-readback}
\lean{EnvelopeProperty.Vanishing.NadelGlobalCanonicalTrace.affineAlgebraCanonicalPullbackMap_raw_apply,EnvelopeProperty.Vanishing.NadelGlobalCanonicalTrace.nativeCanonicalPullbackMap_restricted_toRaw,EnvelopeProperty.Vanishing.NadelGlobalCanonicalTrace.affineAlgebraCanonicalPullbackMap_toRaw_onCarriers,EnvelopeProperty.Vanishing.NadelGlobalCanonicalTrace.affineAlgebraCanonicalPullbackMap_toRaw}
\leanok
Let $k$ be a field and $R,S$ arbitrary commutative $k$-algebras. Let
$\varphi:R\to S$ be the original $k$-algebra homomorphism,
$\beta:\operatorname{Spec}S\to\operatorname{Spec}R$ its spectrum map,
and $f_R,f_S$ the original structural maps to $\operatorname{Spec}k$.
Assume $f_R$ is smooth of relative dimension $n$. Let
$h:f_R\circ\beta=f_S$ be the canonical equality supplied by the
same algebra homomorphism, and let
\[
 \delta:\Omega^n_{f_R}\longrightarrow\beta_*\Omega^n_{f_R\circ\beta}
\]
be the actual adjoint exterior differential. For every original native
affine canonical section $w$ over $R$, reading its original native
algebra-map pullback over $S$ gives exactly
\[
 \operatorname{read}_S\bigl(P_\varphi(w)\bigr)
 =\Omega^n(h)(\operatorname{Spec}S)
   \Bigl(\operatorname{read}_{\beta_*}
          \bigl(\delta(\operatorname{Spec}R)(\operatorname{read}_R(w))\bigr)
   \Bigr).
\]
Every ring, structural map, native section module, and sheaf component in
this identity is the original one. There is no smoothness or integrality
assumption on $S$.
\end{theorem}
\begin{proof}
\uses{lem:kn-canonical-affine-native-readers,lem:kn-canonical-affine-native-raw-map,lem:kn-canonical-affine-native-transport,lem:kn-global-trace-native-branch-differential,lem:kn-global-trace-native-coordinate-wedge}
By its original definition, $P_\varphi$ is the original native pullback
for $\beta$, followed by the original affine equality transport for $h$.
The native affine pullback reads as the actual adjoint exterior
component $\delta$ on the same original section. Apply the field-structure
transport identity to that value. The restriction-of-scalars reader and
pushforward reader keep the section unchanged, giving the displayed
formula. This proves the identity for every section $w$, without a
coordinate expansion, a trace identity, or a coefficient localization.
A bundled commutative ring has the same original carrier and ring
structure; exposing those fields gives the identical statement for
arbitrary original bundled affine rings.
\end{proof}
